\section{Interpolation}

\begin{definition}
    \textbf{Interpolation:} Given exact values 
    $$y_1 = f(x_1,x_2\dots), y_2 = f(x_i, x_{i+1} \dots) \dots$$
    of an unknonw function $f$ determine an interplant $\tilde{f}$ such that 
    $$\tilde{f}(x_1, x_2, \dots) = y_1$$.
\end{definition}

\begin{remark}
    We assume $f$ to have specific properties (smooth, monotone, etc.) $f \in V$ and choose a Subspace $\tilde{f} \in V_n$ together wiht a basis $b \in V_n$
\end{remark}

\begin{example}
    If we choose $V_n$ as the set of all polynomials of degree $\le n$ we are calculating a polynomial interpolation of $f$.
\end{example}

\begin{remark}
    Different spaces and basis are possible and used for different applications.
\end{remark}

\begin{example}
    Using $n+1$ datapoints we can solve the linear system of equations of the $n+1$ datapoints to get a polynomial of degree $n$.
    $$a_n x_1^n + a_{n-1} x_1^{n-1} \dots a_0 = y_1$$
    $$a_n x_2^n + a_{n-1} x_2^{n-1} \dots a_0 = y_2$$
    $$\dots$$
    $$a_n x_{n+1}^n + a_{n-1} x_{n+1}^{n-1} \dots a_0 = y_{n+1}$$
\end{example}

\begin{definition}
    The matrix $A$ formed by this system of linear equations is called a \textbf{Vandermonde matrix}.
\end{definition}

\begin{remark}
    Solving this takes $O(n^3)$ and might come with unreasonably large rounding error.
\end{remark}

\subsection{Newton}

\begin{definition}
    \textbf{Newton-Basis:} The newton basis $N_n(x)$ is 
    $$\prod_{i=0}^{n-1} (x-x_i)$$
\end{definition}

\begin{theorem}
    Using the newton basis we can continously increase the degree of interpolation without having to solve the whole system of equations at once.
    $$p(x_0) = y_0 \implies p_0(x) = y_0$$
    $$p(x_0, x_1) = p_0, y_1 \implies p_1(x) = p_0(x) + \frac{y_1 - y_0}{x_1 - x_0} (x - x_0)$$
    $$p(x_0, x_1, x_2) = p_1, p_1, y_2 \implies p_2(x) = p_1(x) + \frac{\frac{y_2 - y_1}{x_2 - x_1} - \frac{y_1 - y_0}{x_1 - x_0}}{x_2 - x_0} (x - x_0)(x - x_1)$$
    $$\dots$$
\end{theorem}

\begin{remark}
    This can be interpreted as 
    $$p_i(x) = b_1 + b_2 N_1 + b_3 N_2 + \dots$$
    where $b_i$ is the i-th devided difference definied by the follwing recurrence relation
    $$b_i  = f[x_0, \dots x_i] =  \frac{f[x_1, \dots x_i] - f[x_0, \dots x_{i-1}]}{x_i - x_0}$$
    and $f[x_i] = y_i$ 
\end{remark}

\subsection{Lagrange}

\begin{definition}
    \textbf{Lagrange-polynomials} are polynomials of the form
    $$l_i(x) = \prod_{j=0, j\neq i}^{n} \frac{x - x_j}{x_i - x_j}$$
\end{definition}

\begin{theorem}
    Observe that for such a largrange polynomial $l_i(x)$ we have
    $$l_i(x) = 0 \forall x_j \text{with} j\neq i$$
    and 
    $$l_i(x_i) = 1$$
\end{theorem}

\begin{theorem}
    Because the because the above holds, given a set of datapoints $(x_1, y_1)  \dots (x_n, y_n)$ we can construct a polynomial that fullfilles all of the $n$ conditions
    $$p_n(x) = \sum_{i=0}^{n} y_i l_i(x)$$
\end{theorem}

\subsection{Barycentric Interpolation-formula}

\begin{definition}
    \textbf{Barycentric weights:} We define the barycentric weights $w_i$ for a set of nodes $x_0, \dots, x_n$ as
    $$w_i = \prod_{j=0, j \neq i}^{n} \frac{1}{x_i - x_j}$$
    and the node polynomial $\ell(x)$ as
    $$\ell(x) = \prod_{i=0}^{n} (x - x_i)$$
\end{definition}

\begin{theorem}
    Using these definitions, we can rewrite the Lagrange basis polynomials as $l_i(x) = \ell(x) \frac{w_i}{x - x_i}$. Substituting this into the standard Lagrange polynomial yields the \textbf{first form of the barycentric interpolation formula}:
    $$p_n(x) = \ell(x) \sum_{i=0}^{n} \frac{w_i}{x - x_i} y_i$$
\end{theorem}

\begin{theorem}
    \textbf{Derivation of the second form:} Because an interpolating polynomial reproduces constant functions exactly, interpolating the function $f(x) = 1$ (where $y_i = 1 \forall i$) gives:
    $$1 = \ell(x) \sum_{i=0}^{n} \frac{w_i}{x - x_i}$$
    Dividing the first form of the barycentric formula by this identity eliminates the node polynomial $\ell(x)$ and yields the \textbf{second (or true) barycentric interpolation formula}:
    $$p_n(x) = \frac{\sum_{i=0}^{n} \frac{w_i}{x - x_i} y_i}{\sum_{i=0}^{n} \frac{w_i}{x - x_i}}$$
\end{theorem}

\begin{remark}
    The second form is numerically highly stable. Once the weights $w_i$ are precomputed (which takes $O(n^2)$), evaluating $p_n(x)$ at any new point $x$ takes only $O(n)$ operations. Furthermore, adding a new data point only requires updating the weights, which is much cheaper than recalculating the entire Lagrange polynomial.
\end{remark}

\subsection{Chebyshev}

\begin{remark}
    Using equidistant nodes for polynomial interpolation of high degree often leads to severe oscillations at the edges of the interval. This is known as \textbf{Runge's phenomenon}. To minimize this error, we must choose our interpolation nodes strategically.
\end{remark}

\begin{theorem}
    The error of polynomial interpolation is given by:
    $$f(x) - p_n(x) = \frac{f^{(n+1)}(\xi)}{(n+1)!} \prod_{i=0}^{n} (x - x_i)$$
    for some $\xi$ in the interpolation interval. Since we cannot control the $(n+1)$-th derivative of the unknown function $f$, we must minimize the maximum absolute value of the node polynomial $\prod_{i=0}^{n} (x - x_i)$.
\end{theorem}

\begin{definition}
    \textbf{Chebyshev nodes:} The roots of the Chebyshev polynomials of the first kind are called Chebyshev nodes. On the standard interval $[-1, 1]$, the $n+1$ nodes are defined as:
    $$x_i = \cos\left(\frac{2i + 1}{2n + 2}\pi\right) \quad \text{for } i = 0, 1, \dots, n$$
\end{definition}

\begin{remark}
    If the interpolation interval is $[a, b]$ instead of $[-1, 1]$, the nodes can be shifted and scaled using an affine transformation:
    $$x_i^{[a,b]} = \frac{a+b}{2} + \frac{b-a}{2} x_i$$
\end{remark}

\begin{theorem}
    Using Chebyshev nodes minimizes the maximum norm of the node polynomial on $[-1, 1]$:
    $$\max_{x \in [-1, 1]} \left| \prod_{i=0}^{n} (x - x_i) \right| = \frac{1}{2^n}$$
    This node placement guarantees uniform convergence for sufficiently smooth functions and effectively suppresses Runge's phenomenon by clustering the nodes closer together at the edges of the interval.
\end{theorem}

\subsection{Trigonometric Interpolation}

\begin{definition} \textbf{Hilbert Space $L^2(0,1)$:} The space of square-integrable complex-valued functions on the interval $[0,1]$ is defined as:
$$L^2(0,1) := \left\{ f: (0,1) \rightarrow \mathbb{C} \;\middle|\; \|f\|_{L^2(0,1)} = \sqrt{\langle f, f \rangle_{L^2(0,1)}} < \infty \right\}$$
equipped with the inner product $\langle g, f \rangle_{L^2(0,1)} := \int_0^1 \overline{g(t)} f(t) \, dt$.
\end{definition}

\begin{remark} \textbf{$L^2$-Convergence:} A function $f$ is the $L^2$-limit of a sequence of functions $f_n \in L^2(0,1)$ if:
$$\lim_{n \rightarrow \infty} \|f - f_n\|_{L^2(0,1)} = 0 \iff \lim_{n \rightarrow \infty} \int_0^1 |f(t) - f_n(t)|^2 \, dt = 0$$
Unlike pointwise convergence, $L^2$-convergence implies that the total energy (integrated squared difference) between $f$ and $f_n$ vanishes, allowing isolated point discrepancies.
\end{remark}

\begin{theorem} The complex exponential functions $\phi_k(x) = e^{2\pi i k x}$ for $k \in \mathbb{Z}$ form a complete orthonormal basis in $L^2(0,1)$:
$$\langle \phi_k, \phi_j \rangle_{L^2(0,1)} = \int_0^1 e^{-2\pi i k x} e^{2\pi i j x} \, dx = \begin{cases} 1, & \text{if } k = j \\ 0, & \text{if } k \neq j \end{cases}$$
\end{theorem}

\begin{definition} \textbf{Fourier Series:} Any function $f \in L^2(0,1)$ can be represented by its Fourier coefficients $\hat{f}(k)$ as:
$$f(t) = \sum_{k=-\infty}^{\infty} \hat{f}(k) e^{2\pi i k t}$$
where the Fourier coefficients $\hat{f}(k)$ are given by the orthogonal projection onto the basis functions:
$$\hat{f}(k) = \int_0^1 f(t) e^{-2\pi i k t} \, dt$$
\end{definition}

\begin{remark} \textbf{Intuition \& Physical Meaning:} Analogous to decomposing a 3D vector into orthogonal coordinate axes, $f(t)$ is decomposed into pure periodic waves $e^{2\pi i k t}$ of integer frequency $k$. The magnitude $|\hat{f}(k)|$ represents the amplitude of frequency $k$, while its angle represents phase shift. Parseval's identity guarantees energy conservation:
$$\sum_{k=-\infty}^{\infty} |\hat{f}(k)|^2 = \|f\|_{L^2(0,1)}^2$$
\end{remark}

\begin{remark} \textbf{Smoothness and Decay Rate:} The decay rate of Fourier coefficients depends directly on the smoothness of $f$. If $f^{(n)}$ is integrable ($\int_0^1 |f^{(n)}(t)| \, dt < \infty$), then:
$$\hat{f}(k) = \mathcal{O}(k^{-n})$$
Smoother functions yield faster coefficient decay, leading to rapid convergence of the Fourier series.
\end{remark}

\begin{remark} \textbf{Derivative of Fourier Coefficients:} The Fourier coefficients of the $n$-th derivative $f^{(n)}$ satisfy:
$$\widehat{f^{(n)}}(k) = (2\pi i k)^n \hat{f}(k)$$
This converts linear differential operators in the time domain into simple algebraic multiplications in the frequency domain.
\end{remark}

\begin{remark} \textbf{Gibbs Phenomenon:} At jump discontinuities, the truncated Fourier series overshoots near the edge by approximately $9\%$ of the jump size. Increasing the number of terms shifts the overshoot closer to the discontinuity but does not eliminate its amplitude.
\end{remark}

\begin{definition} \textbf{N-th Root of Unity:} The primitive $N$-th root of unity is defined as:
$$\omega_N := \exp\left(\frac{-2\pi i}{N}\right) = \cos\left(\frac{2\pi}{N}\right) - i \sin\left(\frac{2\pi}{N}\right)$$
The roots of unity satisfy the discrete orthogonality relation:
$$\sum_{k=0}^{N-1} \omega_N^{kj} = \begin{cases} N, & \text{if } j \equiv 0 \pmod N \\ 0, & \text{if } j \not\equiv 0 \pmod N \end{cases}$$
This discrete orthogonality serves as the exact algebraic analogue of the continuous inner product of complex exponentials.
\end{definition}

\begin{definition} \textbf{Discrete Fourier Transform (DFT):} Given $N$ equidistant nodes $x_l = \frac{l}{N}$ for $l = 0, \dots, N-1$, approximating the continuous Fourier coefficients $\hat{f}(k)$ via the trapezoidal rule yields the DFT:
$$\hat{f}(k) \approx \hat{f}_N(k) = \frac{1}{N} \sum_{l=0}^{N-1} f\left(\frac{l}{N}\right) \omega_N^{kl}$$
The transformation mapping samples $y = (f(x_0), \dots, f(x_{N-1}))^T \mapsto c = F_N y$ is performed via the Fourier matrix $F_N = (\omega_N^{ij})_{i,j=0}^{N-1}$.
\end{definition}

\begin{remark} \textbf{Fast Fourier Transform (FFT):} Using a divide-and-conquer strategy (e.g., Cooley-Tukey algorithm), the matrix-vector multiplication $F_N y$ can be evaluated in $\mathcal{O}(N \log N)$ operations instead of naive $\mathcal{O}(N^2)$.
\end{remark}

\begin{remark} \textbf{Aliasing:} Sampling $f$ at discrete nodes $x_l = \frac{l}{N}$ makes frequencies differing by integer multiples of $N$ indistinguishable, since $\omega_N^{(k + mN)l} = \omega_N^{kl}$. To avoid aliasing high frequencies into lower modes, $N$ must satisfy the Nyquist-Shannon sampling threshold relative to the highest frequency present in $f$.
\end{remark}

\begin{definition} \textbf{Trigonometric Interpolating Polynomial \& \texttt{fftshift}:} For $N$ samples $y_l = f(t_l)$ at $t_l = \frac{l}{N}$, the trigonometric polynomial $p_N(t)$ satisfying $p_N(t_l) = f(t_l)$ is expressed in symmetric frequency order:
$$p_N(t) = \sum_{k=-N/2}^{N/2-1} \zeta_k e^{2\pi i k t}$$
Standard FFT outputs store coefficients indexed in non-symmetric order $k \in \{0, \dots, N-1\}$. Applying \texttt{fftshift} rearranges indices from $[0, N-1]$ to symmetric frequency spectrum $[-\frac{N}{2}, \frac{N}{2}-1]$.
\end{definition}

\begin{definition} \textbf{Cyclic Convolution:} The cyclic convolution $c = a \otimes b$ of two $N$-periodic vectors $a, b \in \mathbb{C}^N$ is defined as:
$$c_k = \sum_{n=0}^{N-1} a_n b_{(k-n) \pmod N}$$
\end{definition}

\begin{theorem} \textbf{Discrete Convolution Theorem:} The DFT converts cyclic convolution into pointwise vector multiplication:
$$\mathcal{F}_N(a \otimes b) = \mathcal{F}_N(a) \odot \mathcal{F}_N(b)$$
This allows evaluating cyclic convolution in $\mathcal{O}(N \log N)$ operations via:
$$a \otimes b = \text{IFFT}(\text{FFT}(a) \cdot \text{FFT}(b))$$
Linear convolution of length-$N$ sequences can be evaluated using cyclic convolution by zero-padding both sequences to length at least $2N-1$.
\end{theorem}

\begin{remark} \textbf{Zero-Padding for High-Resolution Evaluation:} To evaluate an interpolating trigonometric polynomial $p_N(t)$ at $M \gg N$ fine grid points:
\begin{enumerate}
    \item Compute $c = \text{FFT}(y)$ of size $N$ and center coefficients using \texttt{fftshift}.
    \item Pad the frequency vector with $M - N$ zeros at the high-frequency boundaries.
    \item Apply \texttt{ifftshift} and compute $\tilde{y} = \frac{M}{N} \text{IFFT}(\tilde{c})$ of size $M$.
\end{enumerate}
This yields exact evaluation of $p_N(t)$ on a fine grid without computing individual exponential sums.
\end{remark}
